\documentclass[11pt]{article}
\usepackage[margin=0.82in]{geometry}
\usepackage{amsmath,amssymb,amsthm,booktabs,float,microtype}
\usepackage[font=small,skip=3pt]{caption}
\usepackage[hidelinks]{hyperref}
\setlength{\parindent}{0pt}
\setlength{\parskip}{3pt plus 1pt minus 1pt}
\setlength{\textfloatsep}{6pt plus 2pt minus 2pt}
\setlength{\intextsep}{6pt plus 2pt minus 2pt}
\newtheorem{proposition}{Proposition}
\newtheorem{lemma}[proposition]{Lemma}
\newcommand{\pos}[1]{\left(#1\right)_+}
\newcommand{\AON}{\mathrm{AON}}
\title{\vspace{-1.7em}\textbf{When Selection Breaks Calibration}\\[-1pt]
\large Constructive Collapse in Learned Boolean-Complexity Bounds}
\author{J. R. Landers}
\date{\vspace{-1.4em}}

\begin{document}
\maketitle
\vspace{-1.6em}

\begin{abstract}
A learned model can identify where a mathematical upper bound is loose while
still failing precisely on the cases selected for the largest improvements.
We isolate this phenomenon in a frozen prospective experiment on five-input
Boolean formula complexity.  A semantic model predicted the residual slack in
a certified AND/OR/NOT construction bound; a held-out maximum residual was
subtracted before any reduction was applied.  At exact cost 14, six of 120
refined endpoints fell below the truth.  All six belong to one exact boundary
regime: the construction endpoint is 16, the true value is 15, and therefore
only one gate is removable.  Among ten matched targeted functions in this
regime, the six largest model scores are exactly the six misses.  We derive
the boundary identity behind this correspondence and show how top-score
selection enriches a rare low-headroom subgroup.  Its prevalence rose from
1/46 in cost-13 high-score calibration cases to 9/78 in the cost-14 targeted
high-score cohort.  The functions are structurally distinct, but share a
discordance: restriction nearly closes the bound while global decision-tree
and algebraic features continue to signal difficulty.  The case provides a
finite, reproducible warning that calibration must match the complete
selection policy, not merely the predictor's coarse score stratum.
\end{abstract}

\section{A sharper failure than ordinary prediction error}

Suppose analysis supplies a certified upper bound $Y(f)\le U(f)$ for an
unknown integer-valued quantity $Y(f)$.  A learned model predicts the
\emph{headroom}
\[
 h(f)=U(f)-Y(f)\ge0,
\]
and proposes to remove part of it.  This is more demanding than ordinary
prediction.  Overpredicting $h$ does not merely increase squared error; it can
turn a valid mathematical bound into a false statement.

Our instance is minimum Boolean tree-formula complexity.  For
$f:\{0,1\}^5\to\{0,1\}$, let $K(f)$ be the minimum number of gates in a
formula over $\AON=\{\mathrm{AND},\mathrm{OR},\mathrm{NOT}\}$, and set
$Y(f)=K(f)+1$.  The shift counts formula leaves under the usual tree
convention and aligns with classical lower and constructive upper bounds
\cite{jukna2012}.  The certified endpoint $U(f)$ is the best of a general
analytic construction, exact four-variable restriction constructions, and
disjoint-variable factorizations.

A histogram gradient-boosted regressor \cite{friedman2001} maps 23 semantic
and construction-aware truth-table features to a headroom score $z(f)$.  A
calibration sample supplies a one-sided correction $q_g$ for each score
stratum $g$.  The proposed reduction is
\begin{equation}
 r_\tau(f)=
 \begin{cases}
   \pos{z(f)-q_{g(f)}},&\pos{z(f)-q_{g(f)}}\ge\tau,\\
   0,&\text{otherwise},
 \end{cases}
 \qquad U_{\rm sel}(f)=U(f)-r_\tau(f),                 \label{eq:policy}
\end{equation}
where $\tau=0.25$ suppresses tiny reductions.  On the held-out cost-13 sample,
$q_g$ was the largest dangerous residual $z-h$ in stratum $g$, enlarged by a
cross-layer buffer.  This resembles the basic logic of split calibration
\cite{vovk2005,lei2018}, though here the maximum rather than a nominal
quantile was used.

The frozen cost-14 experiment seemed favorable in aggregate.  It reduced the
construction endpoint on 99 of 120 functions and removed 0.8795 gate on
average.  Nevertheless, coverage was 95\% rather than the required 99\%.
The six misses have a more exact explanation than ``distribution shift.''

\section{The boundary identity and selection enrichment}

Define the amount by which the refined endpoint falls below the truth,
\[
 v(f)=\pos{Y(f)-U_{\rm sel}(f)}.
\]

\begin{lemma}[Headroom boundary]
For the policy in \eqref{eq:policy},
\begin{equation}
 v(f)=\pos{r_\tau(f)-h(f)}.                              \label{eq:violation}
\end{equation}
Whenever $z(f)-q_{g(f)}\ge\tau$ and $h(f)=m$,
\begin{equation}
 v(f)=\pos{z(f)-(q_{g(f)}+m)}.                           \label{eq:boundary}
\end{equation}
Thus an $m$-headroom case fails exactly when $z(f)>q_{g(f)}+m$.
\end{lemma}

\emph{Proof.} Substituting $U_{\rm sel}=U-r_\tau$ and $h=U-Y$ gives
$Y-U_{\rm sel}=r_\tau-h$, proving \eqref{eq:violation}.  In the active branch
of \eqref{eq:policy}, $r_\tau=z-q_g$; substituting $h=m$ proves
\eqref{eq:boundary}. \qed

The lemma shows why small headroom is a boundary layer.  A score may be quite
reasonable for the population yet become unsafe as soon as it crosses one
sharp threshold.  Top-score selection makes this boundary more dangerous.
Let $H_m=\{f:h(f)=m\}$ and let $S_t=\{f:z(f)\ge t\}$ be a score-based selection
event.

\begin{proposition}[Selection enrichment]
For any population with $\Pr(S_t)>0$,
\begin{equation}
 \Pr(H_m\mid S_t)
 =\frac{\Pr(H_m)\Pr(S_t\mid H_m)}{\Pr(S_t)}.             \label{eq:bayes}
\end{equation}
Consequently a rare low-headroom class can be enriched arbitrarily strongly
by top-score selection whenever its upper score tail is heavier than the
population tail.  Marginal validity within a coarse score stratum does not
imply validity conditional on $S_t$ unless the calibration error is stable
under that selection.
\end{proposition}

The identity is Bayes' rule.  Its implication is the important part: selecting
instances because they promise the largest bound reductions is not an
innocent use of a marginally calibrated model.  It changes the mixture of
headroom types inside the evaluation set.

\section{An exact ten-function boundary layer}

The high-stratum cost-14 correction was
\[
 q_{\rm high}=2.360706.
\]
Every miss had $Y=15$ and $U=16$, hence $h=1$.  Lemma~1 therefore places the
failure threshold at
\[
 z>q_{\rm high}+1=3.360706.                              \tag{6}
\]
Exactly ten targeted functions had $U=16$.  Table~\ref{tab:cases} lists all
ten.  The first six scores exceed (6), and all six fail.  The remaining four
scores lie below the boundary, and all four are covered.  The largest score
produces the largest violation automatically:
$3.761361-3.360706=0.400656$.

\begin{table}[H]
\centering\small
\caption{The matched one-gate-headroom group, sorted by frozen model score.
The cofactor column gives the exact AON gate costs of the unique best pair of
four-input restrictions.}
\begin{tabular}{lcrrrrrc}
\toprule
Truth table & Covered & Ones & $z$ & $r$ & $v$ & Leaves & Cofactors\\
\midrule
\texttt{58760a6e} & no  & 15 & 3.7614 & 1.4007 & .4007 & 15 & $4+7$\\
\texttt{b0aae8c0} & no  & 13 & 3.6240 & 1.2633 & .2633 & 14 & $5+6$\\
\texttt{f15c00a8} & no  & 12 & 3.5883 & 1.2276 & .2276 & 13 & $7+4$\\
\texttt{05e402e8} & no  & 11 & 3.5862 & 1.2255 & .2255 & 13 & $8+3$\\
\texttt{0a172737} & no  & 15 & 3.5271 & 1.1664 & .1664 & 12 & $5+6$\\
\texttt{cdd46774} & no  & 18 & 3.3803 & 1.0196 & .0196 & 13 & $5+6$\\
\midrule
\texttt{7676ec28} & yes & 17 & 3.3238 & .9631 & 0 & 12 & $4+7$\\
\texttt{f8ba1102} & yes & 13 & 3.3210 & .9603 & 0 & 13 & $8+3$\\
\texttt{fb8aba32} & yes & 18 & 3.2809 & .9202 & 0 & 13 & $6+5$\\
\texttt{21252111} & yes &  9 & 2.9984 & .8854 & 0 & 12 & $7+4$\\
\bottomrule
\end{tabular}
\label{tab:cases}
\end{table}

This perfect score ordering could still be a duplicated-orbit artifact.  It
is not.  The ten functions occupy ten distinct NPN equivalence classes under
input permutations, input negations, and output negation.  Each has one unique
best restriction variable.  After ignoring cofactor order, their optimal
four-input cost pairs span $3+8$, $4+7$, and $5+6$; covered and uncovered
cases share these patterns.  None is canalizing, and none has a nontrivial
disjoint-variable AND or OR factorization.  There is no single exceptional
normal form hiding behind the six rows.

The commonality is instead a relationship between two kinds of structure.
Restriction has already produced an endpoint just one gate above exact cost,
but the remaining global semantics still appear difficult.  In a local
one-coordinate perturbation against the four covered controls, the largest
upward-score contrasts are decision-tree leaves (0.2104), linear ANF support
(0.1694), quadratic ANF support (0.0938), and degree-two Fourier signs
(0.0566).  Restriction gain has contrast only 0.0203.  These values are
descriptive, not causal, because the coordinates are correlated.  They do,
however, motivate a name for the pattern: \emph{construction--semantics
discordance}.  A specialized construction has nearly closed the bound while
global difficulty features continue to predict room for improvement.

\section{Why calibration did not see the boundary}

The cost-13 calibration sample contained only one direct analogue among its
46 targeted high-stratum functions: a one-headroom case with score 3.099887.
Its dangerous residual was $3.099887-1=2.099887$.  Adding the previous
cross-layer buffer 0.260819 produced the frozen correction 2.360706.

At cost 14, nine of 78 targeted high-stratum cases had one gate of headroom.
The relevant prevalence therefore changed from
\[
 \frac{1}{46}=2.17\%\qquad\text{to}\qquad
 \frac{9}{78}=11.54\%,
\]
a 5.31-fold enrichment.  More importantly, their scores extended to 3.761361,
so the maximum dangerous residual became 2.761361.  The shortfall
$2.761361-2.360706=0.400656$ is exactly the worst violation.

The model did retain useful information.  All 80 targeted functions were
tightened, compared with 19 of 40 random references, and their mean reduction
was 1.1744 rather than 0.2895.  The ranking found opportunity.  The failure is
that calibration pooled unlike headroom regimes and then selection queried
the most aggressive tail of that pool.

This distinction matters for mathematical learning systems.  A safety margin
calibrated on individual predictions need not survive a downstream policy
that searches many candidates and reports only the largest proposed gains.
The calibration unit should be the complete procedure---candidate generation,
scoring, top-$k$ selection, and refinement---or it should condition on enough
construction information to make residual error stable under selection.

\section{Implications}

The case suggests two changes, neither of which can be validated on these same
six functions.  First, the model should expose the complete restriction
witness profile: both cofactor costs, their imbalance, the number of optimal
restriction variables, and the gap to the second-best restriction.  A scalar
restriction gain did not capture how completely restriction had already
explained the function.

Second, the learned target should match the one-sided risk.  Squared-error
prediction of headroom rewards average accuracy.  A direct model of the
dangerous residual $z-h$, a one-sided quantile objective, or calibration that
replays top-$k$ selection would focus on the event that invalidates the bound.
Any new rule derived here is post-hoc.  It must be frozen before a new sample
is labeled.

The six misses therefore form an interesting mathematical case not because
they share one Boolean normal form, but because they instantiate an exact
failure geometry.  Constructive collapse places them one gate from the truth;
semantic persistence assigns large scores; selection concentrates the upper
tail; and the boundary identity converts that tail into six deterministic
violations.  This is a small finite example with a general lesson: when
learning is used to sharpen mathematics, calibration must follow the query
that will actually be made.

\paragraph{Reproducibility.}
The frozen predictions, model and protocol hashes, NPN audit, restriction
witnesses, matched casebook, perturbation table, and support-shift calculation
are stored with the companion experiment.  The structural report is generated
without refitting the model or changing the prospective result.

\begingroup
\footnotesize
\setlength{\parskip}{0pt}
\begin{thebibliography}{9}
\setlength{\itemsep}{0pt}
\bibitem{jukna2012}
S. Jukna, \emph{Boolean Function Complexity: Advances and Frontiers}.
Springer, 2012.
\bibitem{friedman2001}
J. H. Friedman, ``Greedy function approximation: a gradient boosting
machine,'' \emph{Annals of Statistics}, 29(5):1189--1232, 2001.
\bibitem{vovk2005}
V. Vovk, A. Gammerman, and G. Shafer, \emph{Algorithmic Learning in a Random
World}. Springer, 2005.
\bibitem{lei2018}
J. Lei, M. G'Sell, A. Rinaldo, R. Tibshirani, and L. Wasserman,
``Distribution-free predictive inference for regression,''
\emph{J. Amer. Stat. Assoc.}, 113(523):1094--1111, 2018.
\end{thebibliography}
\endgroup

\end{document}
